% =====================================================================
%  THÈME 1 — Pondération des équations chimiques
%  Niveau MOYEN — Corrigé
% =====================================================================
\documentclass[11pt,a4paper]{article}
\usepackage[utf8]{inputenc}
\usepackage[T1]{fontenc}
\usepackage{lmodern}
\usepackage[french]{babel}
\usepackage{amsmath,amssymb,mathtools}
\usepackage{icomma}
\usepackage[version=4]{mhchem}
\usepackage{siunitx}
\sisetup{output-decimal-marker={,},per-mode=symbol}
\usepackage{xcolor}
\usepackage{tikz}
\usepackage[most]{tcolorbox}
\usepackage{enumitem}
\usepackage{array}
\usepackage{booktabs}
\usepackage{fancyhdr}
\usepackage[a4paper,margin=2.1cm,headheight=26pt,headsep=8pt]{geometry}

\definecolor{primary}{RGB}{142,93,168}
\definecolor{primarydark}{RGB}{82,45,108}
\definecolor{corrcol}{RGB}{176,110,20}

\pagestyle{fancy}\fancyhf{}
\fancyhead[L]{\small\textcolor{primarydark}{\textbf{Fiche 5} — Thème 1 : Pondération}}
\fancyhead[R]{\small\textcolor{corrcol}{\textsc{Corrigé — Moyen}}}
\fancyfoot[C]{\small\thepage}
\renewcommand{\headrulewidth}{0.8pt}
\renewcommand{\headrule}{\hbox to\headwidth{\color{primary}\leaders\hrule height \headrulewidth\hfill}}

\newtcolorbox[auto counter]{cor}[1][]{enhanced,breakable,boxrule=1pt,arc=3pt,
  left=3mm,right=3mm,top=2.4mm,bottom=2.4mm,colback=corrcol!6,colframe=corrcol,
  fonttitle=\bfseries,coltitle=white,
  title={Corrigé~\thetcbcounter\ \ifx\relax#1\relax\relax\else\textmd{--- \itshape #1}\fi}}

\newcommand{\rep}[1]{\par\smallskip\noindent\textbf{\color{corrcol}$\Rightarrow$ #1}}

\setlength{\parindent}{0pt}
\setlength{\parskip}{3pt}

\begin{document}

\begin{center}
{\LARGE\bfseries\color{primarydark}Thème 1 — Pondération}\\[1pt]
{\large\color{corrcol}\textbf{Corrigé — Niveau Moyen}}\\
{\color{primary}\rule{0.6\linewidth}{1pt}}
\end{center}

\begin{cor}[combustion du propane]
Carbone : $3$ C $\to 3\,\ce{CO2}$. Hydrogène : $8$ H $\to 4\,\ce{H2O}$.
Oxygène à droite : $3\times2+4\times1=10$ atomes O $\Rightarrow 5\,\ce{O2}$ :
\[ \boxed{\ce{C3H8 + 5 O2 -> 3 CO2 + 4 H2O}} \]
\emph{Vérif :} C $3=3$ ; H $8=8$ ; O $10=10$. \checkmark
\rep{Somme $=1+5+3+4=\boxed{13}$.}
\end{cor}

\begin{cor}[combustion de l'éthane : gérer une fraction]
\textbf{a)} Carbone : $2\,\ce{CO2}$ ; hydrogène : $3\,\ce{H2O}$. Oxygène à droite :
$2\times2+3=7$ atomes, soit des paquets de 2 : il faut $\tfrac{7}{2}\,\ce{O2}$. On obtient
\[ \ce{C2H6 + 7/2 O2 -> 2 CO2 + 3 H2O}\quad(\text{coefficient fractionnaire } \tfrac{7}{2}). \]
\textbf{b)} On multiplie \emph{toute} l'équation par $2$ :
\[ \boxed{\ce{2 C2H6 + 7 O2 -> 4 CO2 + 6 H2O}} \]
\emph{Vérif :} C $4=4$ ; H $12=12$ ; O $14=14$. \checkmark
\rep{Somme $=2+7+4+6=\boxed{19}$.}
\end{cor}

\begin{cor}[un groupement polyatomique intact]
On traite \ce{SO4} en bloc. Aluminium : \ce{Al2(SO4)3} impose 2 Al $\Rightarrow 2\,\ce{Al(OH)3}$.
Groupe \ce{SO4} : 3 à droite $\Rightarrow 3\,\ce{H2SO4}$. Hydrogène : à gauche
$2\times3+3\times2=12$ H $\Rightarrow 6\,\ce{H2O}$ ($6\times2=12$) :
\[ \boxed{\ce{2 Al(OH)3 + 3 H2SO4 -> Al2(SO4)3 + 6 H2O}} \]
\emph{Vérif :} Al $2=2$ ; S $3=3$ ; H $12=12$ ; O : gauche $2\times3+3\times4=18$, droite
$3\times4+6\times1=18$. \checkmark
\rep{Somme $=2+3+1+6=\boxed{12}$.}
\end{cor}

\begin{cor}[une précipitation]
On traite \ce{NO3} en bloc. Aluminium : $1$. Chlore : \ce{AlCl3} donne 3 Cl $\Rightarrow 3\,\ce{AgCl}$,
d'où 3 Ag $\Rightarrow 3\,\ce{AgNO3}$. Groupe \ce{NO3} : $3$ à gauche $=3$ dans \ce{Al(NO3)3} à droite :
\[ \boxed{\ce{AlCl3 + 3 AgNO3 -> 3 AgCl + Al(NO3)3}} \]
\emph{Vérif :} Al $1=1$ ; Cl $3=3$ ; Ag $3=3$ ; \ce{NO3} $3=3$. \checkmark
\rep{Somme $=1+3+3+1=\boxed{8}$.}
\end{cor}

\begin{cor}[équilibrer les charges : équation ionique]
\textbf{a) Atomes.} L'aluminium passe en \ce{Al^3+} (perd 3 électrons), l'hydrogène sort en \ce{H2}.
On équilibre :
\[ \boxed{\ce{2 Al + 6 H+ -> 2 Al^3+ + 3 H2}} \]
Atomes : Al $2=2$ ; H $6=3\times2=6$. \checkmark

\textbf{b) Charges.}
\begin{center}\small
\begin{tabular}{lcc}
\toprule
 & Réactifs & Produits\\ \midrule
Charge & $6\times(+1)=+6$ & $2\times(+3)+3\times0=+6$\\
\bottomrule
\end{tabular}
\end{center}
La charge totale vaut $+6$ des deux côtés : l'équation ionique est correctement pondérée
(atomes \emph{et} charges).
\rep{Somme $=2+6+2+3=\boxed{13}$.}
\end{cor}

\end{document}
